\section{Sheaves on stacks}
\label{sec:sheaves}

\subsection{Quasi-coherent sheaves}
\label{sec:QC}

To any stack $X$, there is an associated stable $\infty$-category $\QC\left(X\right)$ of
quasi-coherent sheaves.
The homotopy category of this is the usual unbounded quasi-coherent derived category.
By definition, $\QC$ is defined on affines $\Spec\left(A\right)$ to be $\Mod_A$.
The definition is extended to a general stack via right Kan extension:
Every stack $X$ can be written as a colimit of affines, and $\QC\left(X\right)$ is then
defined to be the limit of the corresponding diagram in $\Pre$:
\begin{equation*}
\QC\left(X\right)\ceqq \lim_{\Spec A \to X} \Mod_A \in \Pre
\end{equation*}

\subsubsection{Perfect stacks}

We will occasionally make the mild assumption that our stack is well-behaved in the
following sense.

\begin{defn}[{\cite[\S 3.1]{BZFN}}]
\label{defn:perfect}
A stack $X$ is \emph{perfect}
if it has affine diagonal and
$\QC\left(X\right)\simeq \Ind\Perf\left(X\right)$.
\end{defn}

This is a very broad class of stacks, including quasi-compact derived schemes with affine
diagonal, and linear quotients of quasi-projective derived schemes in characteristic zero
\cite[\S 3.3]{BZFN}. In particular, many instances of classifying stacks are perfect.

\begin{theorem}[{\cite[\S 4.2-4.3]{BZFN}}]
\label{thm:BZFN}
Consider maps $X\xrightarrow{p} Y \xleftarrow{p'} X'$ of perfect stacks.
Then $\QC\left(X\right)$ (and $\QC\left(X'\right)$) are self-dual over
$\QC\left(Y\right)$, and there are equivalences
\begin{equation*}
\QC\left(X\right) \otimes_{\QC\left(Y\right)} \QC\left(X'\right) \simeq
\QC\left(X\times _Y X'\right) \simeq
\Fun_{\QC\left(Y\right)}\left(\QC\left(X\right) , \QC\left(X'\right)\right)
\end{equation*}
of $\QC\left(X\right)$-$\QC\left(X'\right)$-bimodule categories.
\end{theorem}

\begin{rmk}
\label{rmk:BZFN_monoidal}
The former equivalence is symmetric monoidal, i.e. an equivalence in
$\CAlg\left(\Pre\right)$.

In the case $X = X'$, the functors from $\QC\left(X\right)$ to itself naturally form a
monoidal category with respect to composition. When considered with this monoidal
structure, the latter equivalence also turns out
to be monoidal, but where the monoidal structure on $\QC\left(X\times_Y X\right)$ is taken
to be convolution (see Section \ref{sec:convolution}).
Explicitly, the latter equivalence sends a kernel $K\in \QC\left(X\times_Y X'\right)$ to
the integral transform $\pi'_{*} \left( \pi^* \left(-\right)\otimes K\right)$, where
$\pi\colon X\times_Y X'\to X$ and $\pi'\colon X\times_Y X'\to X'$ are the projections.
\end{rmk}

\subsubsection{Symmetric-monoidal structure}

For a fixed perfect stack $S$, let $\Stk_S^{\perf}$ be the $\infty$-category of perfect
stacks over $S$, and write
$\Mod_{\QC\left(S\right)}\ceqq\Mod_{\QC\left(S\right)}\left(\Pre\right)$.
Pullback along $X\to S$ makes $\QC\left(X\right)$ a commutative algebra in
$\Mod_{\QC\left(S\right)}$, so we obtain a functor
\begin{equation}
\QC\colon\left(\Stk_S^{\perf}\right)^{\op}\longrightarrow\CAlg\left(\Mod_{\QC\left(S\right)}\right)
\ .
\label{eqn:QC_CAlg}
\end{equation}
Pushouts in $\CAlg\left(\Mod_{\QC\left(S\right)}\right)$ are relative tensor products
\cite[Proposition 3.2.4.7 and Section 4.5]{L:HA}, and $\Stk_S^{\perf}$ is closed under fiber products
\cite[\S 3]{BZFN}.
So Theorem \ref{thm:BZFN} says precisely that \eqref{eqn:QC_CAlg} sends finite limits to
finite colimits.

\subsection{Factorization homology}
\label{sec:fac_hom}

\subsubsection{Viewing the commutative object in the higher Morita categories}
\label{sec:QC_Algn}

Pointwise tensor product of sheaves endows $\QC\left(X\right)$ with a symmetric monoidal structure, so
a stack $X$ defines an object
$\left(\QC\left(X\right) , \otimes\right)\in\umor_d\left(\Pre\right)$ for all $d$.
Its factorization homology is computed by mapping stacks.

\begin{theorem}[{\cite[\S 4]{BZFN}; see \cite{Ste:Tan} for weaker hypotheses}]
\label{thm:int_QC}
Let $X$ be a perfect stack, and let $M$ be a smooth $m$-manifold, $m\leq d$, which is the
interior of a compact manifold with boundary (so that $M$ has finite homotopy type). Since
$\left(\QC\left(X\right),\otimes\right)$ is commutative, its factorization homology over $M$ is
defined without a framing, and agrees with factorization homology of the $\EE_d$-algebra
$\QC\left(X\right)$ over $M\times\RR^{d-m}$. There is a natural symmetric monoidal
equivalence
\begin{equation}
\int_M \QC\left(X\right) \simeq \QC\left(X^M\right) \ ,
\qquad X^M\ceqq\Maps\left(M_B , X\right) \ ,
\label{eqn:int_QC}
\end{equation}
where $M_B$ is the constant (Betti) stack on $M$ and $\Maps$ is the derived mapping stack.
\end{theorem}

\begin{proof}
Since $\left(\QC\left(X\right),\otimes\right)$ is commutative, $\int_M\QC\left(X\right)$ only
depends on the homotopy type of $M$: It is the tensor $M\otimes\QC\left(X\right)$ in
$\CAlg\left(\Pre\right)$, i.e.\ the colimit of the constant diagram with value
$\QC\left(X\right)$ indexed by $M$ \cite[Section 5.5]{L:HA}.
As $M$ is a finite space, this colimit is built from copies of $\QC\left(X\right)$ by finitely
many pushouts, which in $\CAlg\left(\Pre\right)$ are relative tensor products.
Dually, $X^M$ is the corresponding finite limit of copies of $X$, which is again perfect
\cite[\S 3]{BZFN}.
By Theorem \ref{thm:BZFN}, the functor \eqref{eqn:QC_CAlg} sends these finite limits to the
finite colimits computing $M\otimes\QC\left(X\right)$.
\end{proof}

Applying Theorem \ref{thm:BZFN} with $Y = X^2$ (equivalently, Theorem \ref{thm:int_QC} with
$M = S^1$), we obtain the following.

\begin{theorem}[\cite{BZFN}]
\label{thm:BZFN_loops}
For a perfect stack $X$, there are equivalences
\begin{equation*}
\QC\left(X\right) \otimes_{\QC\left(X^2\right)} \QC\left(X\right)
\simeq \QC\left(\mathcal{L}X\right) \simeq
\mathcal{Z}_1\left(\QC\left(X\right) \right) \ .
\end{equation*}
\end{theorem}

\subsection{Dualizability}
\label{sec:QC_dualizability}

\begin{prop}
\label{prop:QC_dualizable}
Let $X$ be a perfect stack. Then $\left(\QC\left(X\right),\otimes\right)$ is
$\left(d+1\right)$-dualizable in $\umor_d\left(\Pre\right)$ for every $d\geq0$.
\end{prop}

\begin{proof}
By Theorem \ref{thm:SSS_dual}, we need $\QC\left(X\right)$ to be dualizable as a module over
\begin{equation*}
\int_{S^{i-1}\times\RR^{d-i+1}}\QC\left(X\right)
\end{equation*}
for $0\leq i\leq d$. Since
$\QC\left(X\right)$ is commutative, this is $\QC\left(X^{S^{i-1}}\right)$ by Theorem
\ref{thm:int_QC}, acting via pullback along the constant maps $X\to X^{S^{i-1}}$ (cf.\ the
proof of Proposition \ref{prop:inv_symm}). The stack $X^{S^{i-1}}$ is a finite limit of copies
of $X$, hence perfect \cite[\S 3]{BZFN}, so $\QC\left(X\right)$ is self-dual over it by Theorem
\ref{thm:BZFN}.
\end{proof}

The argument uses that $\QC\left(X\right)$ is commutative. It does not apply to convolution
categories such as $\left(\QC\left(\cL X\right),*\right)$ or $\DMod^2\left(X\right)$; for the
latter see Corollary \ref{cor:3dual}.

\subsection{The convolution \texorpdfstring{$\EE_2$}{E2}-structure}
\label{sec:LX_convolution}

For a perfect stack $X$, we equip $\QC\left(\cL X\right)$ with the $\EE_2$-monoidal structure
of the Drinfeld center $\mathcal{Z}_1\left(\QC\left(X\right)\right)$, via Theorem
\ref{thm:BZFN_loops}, and write $\left(\QC\left(\cL X\right),*\right)$ for this
$\EE_2$-algebra. Its underlying monoidal structure is convolution relative to $X^2$ (Section
\ref{sec:convolution} with $Y=X^2$): Composition in $\mathcal{Z}_1$ is convolution on $\cL X$.
